\subsubsection{端到端吞吐量与 Batch Size 扩展性}

为评估不同 KV Cache 压缩配置在长上下文解码场景下的端到端性能，本文在 Meta-Llama-3-8B-Instruct 上进行了 batch size 扫描实验。实验设置输入长度为 4096，输出长度为 256，batch size 取 1--6。对比配置包括 Dense、50\% 稀疏、70\% 稀疏、50\% 稀疏 + 2bit 量化以及 70\% 稀疏 + 2bit 量化。

结果表明，两条量化曲线在 batch size 为 1--6 的范围内均稳定超过对应的纯稀疏版本。具体而言，70\% 稀疏配置的吞吐量由 22.35、36.31、53.01、60.21、67.20 和 67.37 tokens/s 提升到 23.10、40.30、53.67、60.41、70.60 和 74.39 tokens/s；50\% 稀疏配置的吞吐量由 22.11、35.16、48.26、56.37、61.45 和 61.80 tokens/s 提升到 22.44、38.32、55.05、62.61、67.91 和 72.03 tokens/s。说明在当前长上下文设置下，优化后的量化稀疏计算核心已经能够将压缩带来的访存收益转化为稳定的端到端吞吐提升。

从不同压缩率的比较结果来看，70\% 稀疏 + 2bit 量化在 batch size 为 1、2、5 和 6 时取得最高吞吐量，而 50\% 稀疏 + 2bit 量化在 batch size 为 3 和 4 时略优。这说明在中等 batch size 区间，更高稀疏度带来的访存收益与量化元数据、反量化开销之间仍存在权衡；但随着 batch size 增大，70\% 稀疏配置的压缩优势会再次占主导。

TTFT 与 TPOT 的结果同样支持这一结论。以 batch size 为 6 为例，70\% 稀疏配置的 TPOT 为 62.59 ms/token，而 70\% 稀疏 + 2bit 量化降低到 56.06 ms/token；50\% 稀疏 + 2bit 量化则为 59.08 ms/token。总体而言，在当前实验设置下，量化版本不仅提升了吞吐量，也改善了中高 batch size 下的稳态解码效率。

需要说明的是，吞吐量基于完整 256-token 生成过程统计；考虑到单卡 24GB 显存限制，TTFT 与 TPOT 采用前 64 个 decode steps 的独立计时窗口估计。图~\ref{fig:end2end-throughput-bs} 与表~\ref{tab:end2end-throughput-bs}--表~\ref{tab:end2end-tpot-bs} 给出了完整结果。
